\section{Two compact hypersurfaces}\label{sec:examples}

Let $X_-$ be the anticanonical threefold for the fan polytope
$\Delta_{25}$ of \cite[\S5]{W5}. For the second example take the
four-dimensional fan polytope whose first four vertices are the
standard basis vectors and whose remaining vertices, in order, are the
columns of
\[
{\setlength{\arraycolsep}{1.2pt}\scriptscriptstyle
\left(\begin{array}{*{23}{r}}
-1&1&1&0&0&-1&-1&1&0&0&-1&0&-1&1&0&-1&1&-2&0&-1&-2&0&-1\\
1&-1&-1&0&0&1&1&-1&0&0&1&-1&0&-2&-1&0&-2&1&-1&0&1&-1&0\\
1&0&-1&1&-1&1&0&-1&0&-1&0&0&1&-1&-1&1&-1&1&0&-1&0&-1&-1\\
0&0&1&-1&1&-1&0&0&-1&0&-1&1&0&1&1&-1&0&-1&-1&0&-1&-1&-1
\end{array}\right).}
\]
This is polytope 528 in the census accompanying \cite{W5}. Number
its vertices from zero. Its four hexagonal faces, written in cyclic
order, are
\[
\begin{aligned}
 F_0&=(19,22,20,17,15,16),& F_1&=(18,15,16,21,24,23),\\
 F_2&=(25,26,23,18,17,20),& F_3&=(25,22,19,21,24,26).
\end{aligned}
\]
Take a general anticanonical hypersurface and subdivide at the
centres of $F_2,F_3$. Write $X_+$ for its proper transform. The two
threefolds $X_-,X_+$ each have thirty ordinary double points and two
anticanonical cones over del Pezzo surfaces of degree six. For $X_+$
this follows by taking the two-faces of the displayed convex hull.
The parallelograms are unit squares up to integral affine
equivalence, the other two-faces apart from the four hexagons are
unimodular triangles, and the dual edge of each parallelogram or
hexagon is primitive. The two centre subdivisions resolve the
corresponding cone points. Resolving the two remaining vertices and
refining projectively gives a common smooth crepant resolution
compatible with the local contractions used below. Its Picard rank is
$27$ for $X_+$ and $23$ for $X_-$. The only nonvertex boundary
lattice points are the hexagon centres, and the dual-edge lengths
give zero correction in the toric divisor count.

\subsection{Supports on the local components}
Write $L$ and $P$ for the line and plane of the local deformation
space. On both threefolds order the two components by $F_0,F_1$, using
the cyclic face markings of \cite[\S5]{W5} on $X_-$. Put the thirty
parallelograms in lexicographic order on their increasing lists of
vertices. The circuit has sign $+1$ at its least vertex and the
opposite vertex. On a marked hexagon let $B_0,\dots,B_5$ denote
the cyclic boundary curves of its vertex resolution. The curve
columns for $L$ are $B_0,B_3$ and those for $P$ are $B_0,B_2,B_4$.
The divisor degrees of $B_j$ are $-1$ at its vertex and the
centre, and $1$ at the two neighboring vertices. The remaining
degrees vanish. These rules give the entire curve matrix.\footnote{The
standalone SageMath supplement \texttt{check\_examples.sage}
reconstructs the faces from the two vertex lists and repeats the
integral calculations of this section.}

In Table~\ref{tab:maximal-support} the node indices run from $0$ to
$29$ and are followed by the two ordered groups of curve indices
over the marked cones. Write $r$ for the rank of all curve classes,
and $r_I$ for the rank of those on a support $I$ of greatest
cardinality. The indicated missing indices vanish in every relation.
The last column records the singularities of a fibre realizing $I$.
\begin{table}[ht]
\centering
\caption{Largest supports for four choices of local components.}
\label{tab:maximal-support}
\begin{tabular}{cccccccl}
\hline
& choice & curves & $r$ & $\dim R$ & $r_I$ & missing indices & singularities\\
\hline
$X_-$&$LP$&35&22&13&21&33&one node\\
$X_+$&$LP$&35&22&13&22&none&smooth\\
$X_+$&$LL$&34&23&11&20&17, 32, 33&node 17 and cone at $F_1$\\
$X_+$&$PP$&36&24&12&21&30, 31, 32&cone at $F_0$\\
\hline
\end{tabular}
\end{table}

Here are integral relations with the asserted supports. The entries
of each vector $d,p,q$ give the thirty old nodal coefficients in
three successive groups of ten.
\[
\begin{array}{c|*{10}{r}}
d&1&-2&1&1&1&-2&-1&1&-1&2\\
 &-2&2&1&1&-1&-1&-1&1&-1&1\\
 &1&1&-1&-2&2&-1&2&1&-2&2\\[2pt]
p&1&2&1&-2&-1&-2&2&-1&2&-1\\
 &1&-1&1&1&-1&2&-2&0&-1&1\\
 &2&-1&2&-1&1&-1&1&-1&-1&1\\[2pt]
q&-2&-1&1&1&-1&-1&2&1&-1&-1\\
 &-2&-2&2&-1&-1&1&1&-2&-1&1\\
 &1&-3&-1&2&-1&2&1&2&2&-2
\end{array}
\]
In table order the four relations are
\[
\begin{aligned}
 &(d,\ 1,-1,\ -1,0,1),&& (p+3q,\ 2,-2,\ 6,-9,3),\\
 &(p,\ 2,-2,\ 0,0),&& (q,\ 0,0,0,\ 2,-3,1).
\end{aligned}
\]
Each relation follows by summing the signed curve columns just
described. Row reduction gives the ranks in the table. Every nonzero
Smith invariant of each active matrix is one.

Write $D(j_1,\ldots,j_t)=\sum_{k=1}^tD_{j_k}$ for the sum of the
indicated prime divisors. The centre rays follow the vertices in face
order. On $X_-$ the divisor $D(2,3,13,14,15,16)$ has degree one on
curve $33$ and degree zero on all other $LP$ columns.
On $X_+$ with choice $LL$, the divisor $D(2,3,4,15,16)$ isolates
the coefficient of node $17$ after the sum-zero relation over
each cone is imposed. The divisor $D(3,6,8,15,17,18,23)$ similarly
isolates the line coefficient at $F_1$. Thus that cone and node
$17$ persist. The two divisor evaluations in the next paragraph
account for the vanished plane at $F_0$ in the last row.

Theorem~\ref{thm:marked-periods} realizes all four displayed
supports. For the two examples with a constant cone, their existence
can also be seen inside the $LP$ deformation space of $X_+$.
Both marked projections are nonzero there. To obtain the $LL$ fibre,
take the line at $F_0$ with coefficients $(2,-2)$ and keep the
plane at $F_1$ at its origin. To obtain the $PP$ fibre, keep the
line at $F_0$ at its origin and take the plane at $F_1$ with
coefficients $(2,-3,1)$. A locally trivial cone deformation belongs
to both components.

\subsection{Divisor classes at the persistent cone}
On each hexagon use the ordered root classes
\[
 \ell=B_0-B_3,\qquad A=B_4-B_0,\qquad B=B_2-B_4 .
\]
They form an integral basis of $K_E^\perp$ in the Picard group of
the degree-six del Pezzo surface. Consider the $PP$ example on
$X_+$. Put
\[
 D_A=-D(11,13,20,22,25,26),\qquad
 D_B=D(3,6,8,15,17,18,23,28).
\]
Both divisors annihilate the thirty old nodes and all three plane
curves over $F_1$. Their degrees on
$(\ell_0,A_0,B_0;\ell_1,A_1,B_1)$, where subscripts now specify
the face, are respectively
\[
 (0,1,0;0,0,0),\qquad (0,0,1;1,0,0).
\]
In particular they isolate the two independent plane coefficients
at $F_0$ when both components are planes. If $N$ denotes the
matrix of the thirty old curve rows, its rank is $20$. Appending
all six root rows increases its rank to $22$. Thus restrictions
of classes annihilating the old curves have rank at most two.
The two tuples above span a primitive sublattice of $\ZZ^6$,
so they give the whole integral image of the surviving classes in
the original marking. On the link of the persistent cone their
degrees span the two plane coordinates. The last three entries in
each tuple retain the marking at the other cone of the central fibre.

The regular locus of this fibre has second Betti number
$27-21-1=5$ by Proposition~\ref{prop:marked-generators}. There are
three further independent classes with zero marked degrees.
Let $H$ be the pullback of the anticanonical polarization on the
original toric fourfold, and let $E_2=D_{29}$ and $E_3=D_{30}$ be the
two exceptional surfaces retained in forming $X_+$. Together
$D_A,D_B,H,E_2,E_3$ give an integral basis modulo torsion. Indeed,
on the resolution one has $H^3=83$ and $H^2D_0=12$, the normalized
volumes of the polar polytope and its facet dual to $v_0$.
The integral class $7HD_0-H^2$ therefore pairs to one with $H$
and pairs to zero with every centre divisor. A boundary line on
each $E_2,E_3$ pairs to $-1$ with its own surface and to zero with
the other surface and $H$. These three functionals annihilate the
class of $E_0$ removed in Proposition~\ref{prop:marked-generators}.
They show that $H,E_2,E_3$ span a primitive sublattice of rank
three in the kernel of the marked degrees. The two displayed root
tuples then give the assertion.
